% STATUT : INSTRUMENT PRÊT, SESSIONS À MENER. C'est le principal apport
% méthodologique restant. Tant que les sessions ne sont pas faites, le
% chapitre tourne sur des experts fictifs, ce qui est un TROU VISIBLE.
\chapter{Élicitation structurée de la propagation}
\label{ch:elicitation}

\textit{[Chapitre à rédiger. Le moteur et le questionnaire existent ; il manque les
séances réelles.]}

% MATÉRIEL DISPONIBLE :
%  - questionnaire_elicitation_w.pdf  : questionnaire AVEUGLE, 10 questions
%    graines + 6 liens cibles, en 5 % / 50 % / 95 %. C'EST CE QUI PART AUX
%    EXPERTS.
%  - protocole_elicitation_w.pdf : la méthode. CONTIENT LES VRAIES VALEURS
%    DES GRAINES. NE JAMAIS L'ENVOYER À UN EXPERT.
%  - antiseche_animation_elicitation.pdf : conduite de séance.
%  - elicitation_reponses_TEMPLATE.csv + script 26b : ingestion et calcul.
%
% PLAN :
%  1. Pourquoi éliciter plutôt que poser : la direction n'est pas dans la
%     donnée (chapitre \ref{ch:identifiabilite}), donc elle vient d'ailleurs.
%     Autant que ce « d'ailleurs » soit traçable et pondéré.
%  2. Le modèle classique de Cooke : questions graines, score de
%     calibration (p-valeur d'un chi2), score d'information (KL contre
%     l'uniforme sur l'étendue COMMUNE aux experts), poids = cal x info
%     avec seuil alpha, décideur = combinaison linéaire pondérée.
%  3. Le dispositif anti-ancrage : un expert à la fois, haute puis basse
%     puis médiane, animateur en miroir neutre.
%  4. Résultats. -- EN ATTENTE DES SÉANCES RÉELLES --
%  5. LIMITE À ÉCRIRE SANS DÉTOUR : panel de petite taille, et certaines
%     graines portent sur des quantités mesurées par le mémoire lui-même.
%     Toute personne ayant vu les résultats empiriques est disqualifiée
%     comme répondant (elle peut relire, pas répondre).
%  6. Ce que l'élicitation alimente exactement : le cran 3 de l'échelle de
%     repli (chapitre \ref{ch:decomposition}), c'est-à-dire u_ij, et rien
%     d'autre. Le reste du modèle est calibré sur données.
